\section{Soundness and Saturation Theorems}\label{sec:soundness}

% Override theorem numbering for this section: v2.1, v2.2 explicit
\renewcommand{\thetheorem}{v2.\arabic{theorem}}
\setcounter{theorem}{0}

We state two theorems. \cref{thm:soundness-v21} (Theorem v2.1)
covers the recovery of states perturbed by an EMR event whose
per-subblock symbol-error count is within the Reed--Solomon
correction capacity. \cref{thm:saturation-v22} (Theorem v2.2)
covers the saturation regime: an event whose per-subblock count
exceeds the capacity yields fail-closed recovery with an auditable
journal entry flagged uncorrectable.

\subsection{Theorem v2.1: recovery under bounded EMR perturbation}

\begin{theorem}[Recovery under bounded EMR perturbation]\label{thm:soundness-v21}
Let~\(\Lambda\) be a layout satisfying invariants
\textbf{I1}--\textbf{I3} of \cref{sub:invariants}, with the INLINE
protection scheme applying \(\mathrm{RS}(n, k)\) per subblock with
\(m\) subblocks per disk block (\(m = 16\) for BEAMFS v2 INLINE). Let~\(\mathcal{G}\) be the recovery
operator of \cref{sec:m2c}, satisfying \textbf{I4}--\textbf{I5}.
Let~\(\mathcal{E}_{\mathrm{EMR}}\) be the family of electromagnetic
perturbation events such that, for every \(e \in \mathcal{E}_{\mathrm{EMR}}\)
and every state~\(s\), the number of corrupted symbols in any single
RS subblock of~\(e(s)\) does not exceed the per-subblock correction
capacity \(t = \lfloor (n-k)/2 \rfloor\) declared by~\(\Lambda\).
Then for every initial consistent state~\(s_0\) (i.e.\ \(s_0 \models\) per \cref{def:consistent}) and every
\(e \in \mathcal{E}_{\mathrm{EMR}}\), there exists~\(K \le |\Lambda|\)
such that
\[
  \mathcal{G}^{K}\bigl(e(s_0)\bigr) \models
  \quad\text{or}\quad
  \mathcal{G}^{K}\bigl(e(s_0)\bigr) = \bot,
\]
where~\(|\Lambda|\) denotes the number of regions in~\(\Lambda\).
\end{theorem}

The full proof appears in \cref{app:proof}. The argument is the
same termination/safety/coverage structure as in the v1 retracted
theorem; what changes is the family quantified over.

\paragraph{Termination.}
By Lemma~\ref{lem:descent} (\cref{app:proof}), the function~\(N\)
counting non-consistent regions strictly decreases on each
non-fail-closed iteration. Since~\(N \in \{0, 1, \dots, |\Lambda|\}\)
and~\(N = 0\) implies consistency by \cref{def:consistent}, the
recursion terminates in at most~\(|\Lambda|\) iterations.

\paragraph{Safety.}
At every iteration, if any step of \(\mathcal{G}\) detects a CRC
mismatch after RS decode or a sentinel violation,~\(\bot\) is
returned. By \cref{prop:absorb} (\cref{app:proof}), fail-closed is
absorbing: no subsequent iteration can return a state that
silently masks the violation.

\paragraph{Coverage of the perturbation family.}
The hypothesis on~\(\mathcal{E}_{\mathrm{EMR}}\) (no event corrupts
more than~\(t\) symbols in any single subblock) is exactly the RS
code's correction capacity. Events that violate this hypothesis
are addressed by \cref{thm:saturation-v22}.

\subsection{Theorem v2.2: saturation observability}

\begin{theorem}[Saturation observability]\label{thm:saturation-v22}
Let~\(\Lambda\) and~\(\mathcal{G}\) be as in
\cref{thm:soundness-v21}. Let~\(\mathcal{E}_{\mathrm{sat}}\) be the
family of electromagnetic perturbation events such that, for some
\(e \in \mathcal{E}_{\mathrm{sat}}\) and some state~\(s\), there
exists at least one RS subblock of~\(e(s)\) whose symbol-error
count exceeds~\(t\). Then for every consistent~\(s_0\) and every
\(e \in \mathcal{E}_{\mathrm{sat}}\):
\begin{enumerate}
\item \(\mathcal{G}\bigl(e(s_0)\bigr) = \bot\); and
\item the persistent journal region of~\(\Lambda\) contains, after
  the application of~\(\mathcal{G}\), an entry tagged with
  \texttt{BEAMFS\_RS\_EVENT\_FLAG\_UNCORRECTABLE} whose
  \((\text{block},\text{subblock})\) coordinates identify a
  saturating subblock of~\(e(s_0)\).
\end{enumerate}
\end{theorem}

The proof appears in \cref{app:proof}. The argument has two parts.

\paragraph{Fail-closed.}
The recovery operator's RS decode step on a saturating subblock
returns a non-zero residual syndrome after decode (the canonical
behaviour of an RS decoder beyond its capacity~\cite{kernelrslib}).
The operator interprets this residual as a saturation event,
returning~\(\bot\) by I5.

\paragraph{Observability.}
Before returning~\(\bot\), the operator emits a journal entry by
calling \texttt{beamfs\_log\_rs\_event} with positions \(=\)
\texttt{NULL} and \texttt{n\_positions} \(=\) 0, accepted by the
on-disk format v4 specification as the canonical encoding of a
saturating event. The entry includes the affected
\((\text{block},\text{subblock})\) coordinates and the
\texttt{BEAMFS\_RS\_EVENT\_FLAG\_UNCORRECTABLE} flag (bit 1 of
\texttt{re\_flags}~\cite{beamfs-impl}). The entry is committed to
persistent storage before the operator returns.

\subsection{Discussion}

\paragraph{What these theorems do, and do not, prove.}
\Cref{thm:soundness-v21,thm:saturation-v22} establish that the
recovery operator either returns to consistency or to fail-closed
in a bounded number of steps under the EMR-bounded hypothesis,
and that fail-closed is auditable, in the sense that a subsequent
inspection of the journal identifies the saturating coordinates.
They do \emph{not} prove that the kernel implementation of
\(\mathcal{G}\) is bug-free: a buggy implementation can violate
\textbf{I4} or \textbf{I5} even with a correct operator
definition. \Cref{sec:m2c} provides a traceability table pairing
each operator step with the kernel function that realizes it; an
implementation audit against this table is a necessary condition
for the theorems to apply to the implementation, not a sufficient
one. A verified compilation chain (in the seL4~\cite{klein2009sel4}
or FSCQ~\cite{chen2015fscq} sense) remains the only way to close
the implementation gap; we mark this as v3+ work.

\paragraph{Why the bound is~\(|\Lambda|\).}
The bound scales with the number of regions because the worst
case is a perturbation that corrupts every region simultaneously
up to the per-subblock capacity. This worst case is unrealistic
for Family A under environmental backgrounds: \cref{sec:emr-threat}
establishes that simultaneous saturation of all regions requires
fluences orders of magnitude beyond environmental SEU rates. For
Family B, simultaneous saturation is the adversary's stated
objective, so the conservative bound is appropriate. The
per-iteration cost is dominated by a single \(\mathrm{RS}(255,239)\)
decode of \(O(n^2)\) GF operations, sub-millisecond on
contemporary CPUs.

\paragraph{Comparison with crash-consistency proofs.}
FSCQ~\cite{chen2015fscq} proves that a filesystem returns to a
consistent state after a crash. The hypothesis there is that the
storage substrate is \emph{benevolent}: it preserves whatever was
written, exactly. The present soundness theorem assumes the
substrate is \emph{adversarial} within the EMR-bounded family,
and the filesystem must recover. The two proofs compose: a
BEAMFS-compliant implementation satisfying
\cref{thm:soundness-v21,thm:saturation-v22} \emph{and} satisfying
crash-Hoare logic in the FSCQ sense recovers from the union of
both threat models. Mechanizing this composition is left for
future work.

\paragraph{Falsifiability.}
By construction, both theorems are falsifiable: a counter-example
is any RadFI~\cite{desbrieres2026radfi} run, or any equivalent
fault-injection trace, that exhibits a state \(e(s_0)\) under
\(e \in \mathcal{E}_{\mathrm{EMR}}\) for which
\(\mathcal{G}^K(e(s_0)) \neq \models\) and
\(\mathcal{G}^K(e(s_0)) \neq \bot\); or a state under
\(e \in \mathcal{E}_{\mathrm{sat}}\) for which \(\mathcal{G}\) does
not return \(\bot\) or fails to emit the
\texttt{UNCORRECTABLE}-flagged journal entry. The empirical
section (\cref{sec:empirical}) reports a corroboration of
\cref{thm:soundness-v21} by direct disk corruption; saturation
testing (Theorem v2.2 active falsification) is part of the
post-Zenodo continuation.
